\begin{example}[A sheaf is fixed by sheafification]
Let $\mathcal F=C^0_X$, the sheaf of continuous real-valued functions.  Every
$\sigma\in\mathcal F^\#(U)$ is locally represented by ordinary continuous functions.  These
local representatives agree on overlaps because they determine the same germs, hence glue
inside $C^0_X$.  Thus every sheafified section comes from one section of $C^0_X(U)$ and
\[
C^0_X\xrightarrow{\sim}(C^0_X)^\#.
\]
This is the concrete version of idempotence.
\end{example}

\begin{example}[Constant presheaf on a disconnected open]
Let $A$ be an abelian group and let $\mathcal P(U)=A$ for every nonempty open $U$, with
identity restriction maps.  Suppose $U=U_1\sqcup U_2$ with both pieces nonempty.  A section
of $\mathcal P(U)$ is one element $a\in A$, so its restrictions have the same value on both
components.  In the sheafification, however,
\[
\mathcal P^\#(U)=\{\text{locally constant maps }U\to A\}\cong A\times A,
\]
and $(a,b)$ with $a\neq b$ is a new section added by sheafification.
\end{example}

\begin{example}[Constant presheaf on a connected interval]
On an interval $I\subseteq\RR$, every locally constant $A$-valued function is constant.
Hence
\[
\mathcal P^\#(I)\cong A\cong\mathcal P(I).
\]
The constant presheaf can therefore look correct on each connected open while still failing
on disconnected opens.  Sheafification repairs the global bookkeeping without changing any
stalk, all of which are $A$.
\end{example}

\begin{example}[Bounded continuous functions]
Let $\mathcal B(U)$ be the bounded continuous functions $U\to\RR$.  VI/13 showed that this
is not a sheaf.  Every continuous function is locally bounded: if $f(x)=c$, continuity gives
an open neighborhood on which $f$ takes values in $(c-1,c+1)$.  Therefore every section of
$C^0_X$ is locally represented by a section of $\mathcal B$, and conversely bounded
continuous functions are continuous.  Hence
\[
\mathcal B^\#\cong C^0_X.
\]
Sheafification removes the nonlocal condition ``bounded on all of $U$.''
\end{example}

\begin{example}[Locally polynomial functions on the real line]
Define $\mathcal P(U)$ to be restrictions to $U$ of a single polynomial in $\RR[t]$.  This
is a presheaf.  On a disconnected open $U=U_1\sqcup U_2$, one may choose different
polynomials on the two components, so $\mathcal P$ is generally not a sheaf.  Its
sheafification consists of functions which are locally polynomial.  On each connected open
interval such a function is represented by one polynomial, since two polynomials agreeing
on a nonempty open interval are equal.
\end{example}

\begin{example}[Two sections become equal after sheafification]
Let $X$ be a space in which every point has a proper open neighborhood.  Define a presheaf
of abelian groups by
\[
\mathcal F(X)=A,\qquad \mathcal F(U)=0\quad(U\subsetneq X),
\]
with the only possible restriction maps.  Every germ of every global section is zero,
because near each point one can restrict to a proper open.  Consequently every stalk is
zero and
\[
\mathcal F^\#=0.
\]
Distinct elements of $A=\mathcal F(X)$ are all identified by sheafification.
\end{example}

\begin{example}[A presheaf with the same stalks as the zero sheaf]
In the preceding example, $\mathcal F_x=0$ for every $x$, even though $\mathcal F(X)=A$ can be nonzero.  The canonical map
\[
\mathcal F\to0
\]
is a stalkwise isomorphism and induces an isomorphism after sheafification.  This shows that
stalks are local invariants, not a complete invariant of arbitrary presheaves.
\end{example}

\begin{example}[One-point space]
Let $X=\{x\}$.  The only opens are $\varnothing$ and $X$.  A presheaf of abelian groups with
$\mathcal F(\varnothing)=0$ is already a sheaf: every nonempty open cover of $X$ contains
$X$.  Therefore
\[
\mathcal F^\#\cong\mathcal F.
\]
The espace \'{e}tal\'{e} is simply the discrete fiber $\mathcal F_x=\mathcal F(X)$ over the
single point.
\end{example}

\begin{example}[Two-point discrete space]
Let $X=\{p,q\}$ be discrete and start from the naive constant presheaf $A$.  Its stalks are
$A$ at both points, so the sheafification is the sheaf determined by these two stalks:
\[
\mathcal F^\#(X)=A\times A,
\qquad
\mathcal F^\#(\{p\})=A,
\qquad
\mathcal F^\#(\{q\})=A.
\]
The unit on $X$ is the diagonal map $A\to A\times A$.
\end{example}

\begin{example}[A Sierpi\'{n}ski space]
Let $X=\{\eta,x\}$ with opens $\varnothing,\{\eta\},X$.  A presheaf of abelian groups is a
map $A=\mathcal F(X)\to B=\mathcal F(\{\eta\})$.  Every cover of $X$ contains $X$ itself,
so every such presheaf is already a sheaf.  Its sheafification is therefore unchanged.  This
finite example emphasizes that failure of the sheaf axiom depends on the available covers.
\end{example}

\begin{example}[The canonical map on stalks]
Let $s\in\mathcal F(U)$ and $x\in U$.  Under
\[
\mathcal F_x\xrightarrow{\sim}(\mathcal F^\#)_x,
\]
the germ $s_x$ is sent to the germ $(\widetilde s)_x$.  The inverse evaluates a sheafified
germ at $x$.  Thus the stalk isomorphism is not merely existential; it is the obvious map
suggested by the espace \'{e}tal\'{e} construction.
\end{example}

\begin{example}[Local factorization into a sheaf]
Let $u:\mathcal F\to C^0_X$ be a presheaf morphism and
$\sigma\in\mathcal F^\#(U)$.  Choose a cover $U=\bigcup U_i$ with
$\sigma|_{U_i}=\widetilde{s_i}$.  The sections $u(s_i)$ of $C^0_X(U_i)$ have equal germs on
overlaps, hence are equal there and glue to one continuous function on $U$.  This is exactly
$\bar u_U(\sigma)$ in the universal property.
\end{example}

\begin{example}[Naturality of the unit]
For a morphism $f:\mathcal F\to\mathcal G$ and a section $s\in\mathcal F(U)$, the two routes
around the naturality square give
\[
f^\#(\widetilde s)=\widetilde{f(s)}.
\]
Indeed, at $x$ both sides have value $f_x(s_x)$.  Thus the commutative square for the unit is
visible point by point in the fibers of the two espace \'{e}tal\'{e}s.
\end{example}

\begin{example}[Sheafification of the naive constant presheaf on a circle]
Let $A=\ZZ$ and $X=S^1$.  The sheafification is the locally constant integer-valued sheaf.
Since $S^1$ is connected,
\[
\underline\ZZ(S^1)=\ZZ.
\]
But on a disconnected open consisting of two short arcs,
\[
\underline\ZZ(U)\cong\ZZ\times\ZZ.
\]
Thus global sections on the whole circle happen to agree with the naive presheaf, while the
presheaf still fails to be a sheaf on smaller disconnected opens.
\end{example}

\begin{example}[Sheafification commutes with open restriction]
Let $j:U\hookrightarrow X$.  Both $(\mathcal F|_U)^\#$ and
$\mathcal F^\#|_U$ have the same espace \'{e}tal\'{e} over $U$:
\[
\coprod_{x\in U}\mathcal F_x.
\]
Their locally representable sections are therefore the same.  Hence
\[
(\mathcal F|_U)^\#\cong\mathcal F^\#|_U.
\]
\end{example}

\begin{example}[A basis computation]
Let $X=\RR$ and let $\mathcal B$ be the basis of open intervals.  Suppose a presheaf is given
only on intervals and satisfies compatible restriction maps.  The germ at $x$ can be
computed using only intervals containing $x$, because they are cofinal among neighborhoods.
The espace \'{e}tal\'{e} construction then produces the same sheaf as if one first extended
the presheaf to all opens and then sheafified.
\end{example}

\begin{example}[The first plus construction removes ambiguity]
Suppose $s,t\in\mathcal F(U)$ restrict equally on an open cover $U=\bigcup U_i$ although
$s\neq t$.  In the plus construction, both sections determine the same compatible family
$\{s|_{U_i}\}=\{t|_{U_i}\}$, so their images in $\mathcal F^+(U)$ coincide.  Thus
$\mathcal F^+$ satisfies locality even if $\mathcal F$ does not.
\end{example}

\begin{example}[The second plus construction supplies gluings]
Suppose $\mathcal F$ is already separated but compatible sections on a cover fail to glue.
Their compatible family itself defines an element of $\mathcal F^+(U)$.  Because separatedness
removes ambiguity on refinements, these new local objects glue correctly; standard plus
construction theory shows $\mathcal F^+$ is now a sheaf.  For arbitrary $\mathcal F$, two
applications give $\mathcal F^{++}\cong\mathcal F^\#$.
\end{example}

\begin{example}[Idempotence through the universal property]
Let $\mathcal G$ be a sheaf.  The identity map $\mathcal G\to\mathcal G$ factors uniquely
through $\eta_{\mathcal G}:\mathcal G\to\mathcal G^\#$.  The resulting map
$r:\mathcal G^\#\to\mathcal G$ satisfies $r\eta=\id$.  Applying uniqueness again gives
$\eta r=\id$.  This constructs the canonical inverse to the unit without inspecting any
sections explicitly.
\end{example}

\begin{example}[Same sheafification, different presheaves]
On a disconnected space, let $\mathcal P$ be the naive constant presheaf $A$ and let
$\underline A$ be the constant sheaf.  They are not isomorphic as presheaves, because on a
disconnected open $U$ the first has $A$ while the second can have a product of several
copies of $A$.  Nevertheless
\[
\mathcal P^\#\cong\underline A\cong(\underline A)^\#.
\]
Sheafification forgets the defective global presentation and retains the same local germs.
\end{example}
